\documentclass[10pt,a4paper]{article}
\usepackage[margin=1.9cm]{geometry}
\usepackage[T1]{fontenc}
\usepackage{lmodern}
\usepackage{amsmath,amssymb}
\usepackage{booktabs,tabularx}
\usepackage{enumitem}
\usepackage[dvipsnames]{xcolor}
\usepackage[most]{tcolorbox}
\usepackage[colorlinks=true,linkcolor=blue!60!black,urlcolor=blue!60!black,citecolor=blue!60!black]{hyperref}
\setlist{nosep,leftmargin=1.2em}
\setlength{\parindent}{0pt}
\setlength{\parskip}{4pt}
\newcommand{\code}[1]{\texttt{\detokenize{#1}}}
\newtcolorbox{keybox}[1]{colback=blue!3,colframe=blue!45!black,fonttitle=\bfseries,title=#1,boxrule=0.6pt,arc=1mm,left=4pt,right=4pt,top=2pt,bottom=2pt}

\title{\textbf{Pole placement and feedforward fit for the CNC plotter}\\[2pt]\large what was done, why, and on what it rests}
\date{}
\begin{document}
\maketitle
\vspace{-9mm}

\begin{keybox}{The workflow in five sentences}
\begin{enumerate}
  \item The joint is modelled as $G(s)=b/(s(\tau s+1))$, identified on the robot; an integral state is added and the three
        closed-loop poles are placed (\code{design_pp.m}), with a Luenberger observer for the speed.
  \item The poles follow one pattern (complex pair $\omega_n,\zeta$ + real integral pole $\lambda$), so every parameter has a
        physical meaning.
  \item The pattern is tuned by a constrained search: the poles move as little as possible from a fast starting design until
        every robustness and actuator requirement is met, \emph{including the measured delay}.
  \item The feedforward (inverse model) is then corrected from a run: a linear least-squares fit of the logged voltage on the
        reference speed, acceleration and friction. It is outside the loop, so it does not change the design.
  \item Result: the PP meets $PM\geq45^\circ$, $M_s\leq1.6$ with the delay, and with the fitted feedforward the tracking error
        returns to the level without disturbances, below one encoder count at the tip.
\end{enumerate}
\end{keybox}

%==========================================================================================
\section{Pole placement: the design steps and their justification}

\textbf{1. Model and feasibility.} State $x=[\theta,\omega]^T$ plus the integral of the error $e_I$. Pole placement with arbitrary
poles exists if and only if the pair is reachable; an asymptotic observer with arbitrary poles exists if and only if $(A,C)$ is
observable (course notes, \emph{Pole placement}; Ackermann's formula~\cite{ackermann72}). Both hold here
(the script checks the ranks). The integral state gives zero steady-state error to constant references and load
disturbances (internal model principle).

\textbf{2. Observer.} Only $\theta$ is measured, so $\omega$ is estimated by a Luenberger observer~\cite{luenberger71}. By the
separation principle the closed-loop eigenvalues are the union of those of $A-BK$ and $A-LC$ (course notes), so controller and
observer are designed separately. As in the notes, the observer poles are chosen \emph{much faster} than the controller ones
(6 times, capped at $250$ rad/s by the 2 ms sampling); being an algorithm, the observer does not saturate.

\textbf{3. Where to put the poles.} The course notes suggest imposing an approximately first/second-order behaviour with a
dominant pole pair, and warn against a closed loop much faster than the open loop, because of actuator saturation. Design A
does exactly this:
\begin{equation}
  p_{1,2}=-\zeta\omega_n\pm j\omega_n\sqrt{1-\zeta^2},\qquad p_3=-\lambda ,
\end{equation}
with the voltage peak and the noise on the voltage among the requirements, so that the speed is limited by the actuator.

\textbf{4. Requirements} (\code{pp_setup.m}), each tied to a physical reason:

\begin{tabularx}{\linewidth}{@{}l c X@{}}
\toprule
requirement & value & reason \\
\midrule
peak voltage & 6 V & saturation of the amplifier: the linear design stays valid \\
$M_s=\|1/(1+L)\|_\infty$ & $\leq1.6$ & distance of the Nyquist curve from $-1$; implies $GM\geq M_s/(M_s-1)=8.5$ dB and
  $PM\geq2\arcsin(1/2M_s)=36^\circ$~\cite{skogestad05}; typical robust range $1.4$--$2$~\cite{astrom06} \\
phase margin & $\geq45^\circ$ & robustness to delay and to errors in $k,\tau$ \\
noise on $u$ & $\leq0.5$ V & one encoder count times the controller gain at high frequency \\
slowest pole & $\lambda\geq10$ rad/s & the integral must recover the friction error within one side of the square (1 s) \\
delay & $T_d=1.5T_s+3.5$ ms & sample-and-hold + computation + measured input delay; a delay removes phase $\omega T_d$ \\
\bottomrule
\end{tabularx}

\textbf{5. Reference weight $\beta=1$} for trajectories: $\beta$ only moves the zero of the reference-to-angle transfer
(set-point weighting~\cite{astrom06}); with $\beta=1$ the feedback acts on $\theta-\phi$ along the whole path, so it does not lag.
$\beta^*$, which cancels the real pole, is kept for point-to-point steps.

\textbf{6. Anti-windup}: back-calculation, $\dot e_I=e+K_{aw}(u_{sat}-u)$~\cite{astrom06}: when the amplifier saturates, the
integral is unloaded instead of accumulating.

%==========================================================================================
\section{The search: why this cost function}

The three parameters $x=[\omega_n,\zeta,\alpha]$ (with $\lambda=\alpha\zeta\omega_n$) are found by
\begin{equation}
  \min_x\;\underbrace{\sum_{i=1}^3\bigl(\log|p_i(x)|-\log|p_{i,0}|\bigr)^2}_{\text{stay close to the designer's start}}
  \;+\;w\underbrace{\sum_k\max(0,g_k(x))^2}_{\text{violated requirements}},\qquad w=10^4 .
\end{equation}
\begin{itemize}
  \item \textbf{Exterior quadratic penalty}~\cite{nocedal06}: a constrained problem is turned into an unconstrained one;
        the penalty is zero inside the feasible set, so among feasible designs the closest to the start is chosen. A large $w$
        means: first feasibility, then closeness.
  \item \textbf{Normalised violations} $g_k=\text{value}/\text{limit}-1$: requirements in volts, degrees and pure numbers weigh
        the same.
  \item \textbf{Logarithms}: distances between poles are relative (a 10\% move counts the same at 5 or 50 rad/s), and the
        parameters stay positive without extra constraints.
  \item \textbf{Nelder--Mead} (\code{fminsearch})~\cite{nelder65,lagarias98}: the requirements come from simulations and
        frequency responses, so the cost has no analytical gradient and is not smooth; a derivative-free simplex method in
        three dimensions is the standard choice. It finds a local minimum, which is why the start is chosen by the designer.
  \item The same idea, optimising a performance index under an $M_s$ constraint, is the basis of robust PI/PID design
        methods~\cite{astrom98,panagopoulos02,garpinger08}.
\end{itemize}

\textbf{What the search did.} From $\omega_n=45$, $\zeta=0.9$, $\lambda=12$ (fast on purpose) it moved to
$\omega_n=31.8$, $\zeta=0.999$, $\lambda=10$. With the old 3 ms delay the start looked feasible ($PM=49^\circ$, $M_s=1.53$);
with the measured 6.5 ms it was not ($PM=39^\circ$, $M_s=1.96$). The final design has $PM=50^\circ$, $M_s=1.59$ on the model
and $PM\approx45^\circ$, $M_s\approx1.72$ on the plant with $+23\%$ gain: the delay, not the controller structure, sets the
achievable speed.

%==========================================================================================
\section{The feedforward fit: why it is sound}

\textbf{Why a feedforward.} The tracking error is $e=S\,(1-PF)\,\phi$, with $S=1/(1+PC)$: the feedforward $F$ removes the error
when $F=P^{-1}$, and it does not enter $S$, so it cannot affect stability~\cite{lambrechts05,boerlage03}. For the motor,
$P^{-1}$ is a speed term and an acceleration term (rigid-body feedforward~\cite{lambrechts05}), plus a Coulomb friction
compensation~\cite{armstrong94}:
\begin{equation}
  u_{ff}=-\frac{1}{k}\dot\phi-\frac{\tau}{k}\ddot\phi-V_c\tanh\!\left(\frac{\dot\phi}{w_\varepsilon}\right).
\end{equation}

\textbf{Why fit it on data.} The identified $k,\tau$ are not exact, and friction and delay are not in the linear model. In the
periodic regime the angle follows the path, so the logged voltage is what the real motor needs. Since $u_{ff}$ is
\emph{linear in its parameters}, fitting it to the logged voltage is a linear least-squares problem with a unique solution
\cite{ljung99}. This is the same principle as fixed-structure feedforward tuning on the real machine, where a linear
parametrisation gives a convex problem whose optimum is with respect to the actual plant~\cite{vandermeulen08}; here it is
solved in one shot instead of iteratively, which is possible because the voltage is logged.

\textbf{Details that make it correct.}
\begin{itemize}
  \item Regressors from the \emph{reference} (smooth, noise free), not from the measured angle: no amplified quantisation and no
        errors-in-variables bias~\cite{soderstrom07}.
  \item Only the periodic regime is used, without saturated samples or disturbance tests.
  \item A small delay is absorbed as $\tau_{ff}\approx\tau+T_d$ (first-order Taylor expansion), so it is compensated without
        modelling it.
  \item Only $\phi$ and $u$ are used: the same function works for the PID; a run with any controller is valid data.
\end{itemize}

\textbf{Workflow.} Simulation: one run with the model feedforward, fit, run with the fitted one (automatic with
\code{ff_source='sim'}). Robot: one run (or an existing one with the same path) with the model feedforward, fit on its data,
runs with the fitted feedforward. This is the standard practice of tuning the feedforward on the machine.

%==========================================================================================
\section{Results and their limit}

\begin{itemize}
  \item With delay and friction, PP + model feedforward has a clearly larger error than without; with the fitted feedforward it
        returns to the level without disturbances. The PID under the same conditions has a larger error.
  \item The contour error (distance from the path, the quality of the drawing~\cite{koren80}) is close to zero; the remaining
        tracking error ($\approx0.16$ mm) is a small timing lag along the path.
  \item $0.16$ mm is below one encoder count at the tip ($0.127\,\text{m}\cdot2\pi/4096\approx0.19$ mm): the sensor cannot
        resolve it, so no controller can reduce it further with this measurement. What remains comes from the coupling between the
        links, which a single-joint model does not contain; LQG (noise) and MPC with a coupled model address different parts of
        this.
\end{itemize}

\section*{Likely questions}
\begin{tabularx}{\linewidth}{@{}>{\itshape}p{5.2cm} X@{}}
\toprule
Why not place the poles by hand? & The requirements interact (speed vs.\ $M_s$ vs.\ noise); the search makes the trade-off
  explicit and reproducible, starting from a hand-chosen design. Hand tuning tools exist and give the same result. \\
Why is $\zeta=0.999$? & With a delay, a well damped pair gives the lowest $M_s$ for a given speed; the pattern is kept complex by
  capping $\zeta<1$. \\
Is the feedforward fit identification in closed loop? & Yes, but of the \emph{inverse} plant along a known reference; the
  regressors are noise free and independent of the controller, so the usual closed-loop bias does not arise. \\
Does the fit change stability? & No: the feedforward is outside the loop ($F$ does not appear in $S$). \\
Why not drive the error to zero? & It is below the encoder resolution and comes from unmodelled coupling. \\
\bottomrule
\end{tabularx}

\begin{thebibliography}{99}\small
\setlength{\itemsep}{1pt}
\bibitem{ackermann72} J.~Ackermann, ``Der Entwurf linearer Regelungssysteme im Zustandsraum,'' \emph{Regelungstechnik}, vol.~20, pp.~297--300, 1972.
\bibitem{luenberger71} D.~G.~Luenberger, ``An introduction to observers,'' \emph{IEEE Trans.\ Automatic Control}, vol.~16, no.~6, pp.~596--602, 1971.
\bibitem{skogestad05} S.~Skogestad and I.~Postlethwaite, \emph{Multivariable Feedback Control: Analysis and Design}, 2nd ed., Wiley, 2005.
\bibitem{astrom06} K.~J.~\AA{}str\"om and T.~H\"agglund, \emph{Advanced PID Control}, ISA, 2006.
\bibitem{nocedal06} J.~Nocedal and S.~J.~Wright, \emph{Numerical Optimization}, 2nd ed., Springer, 2006 (ch.~17, penalty methods).
\bibitem{nelder65} J.~A.~Nelder and R.~Mead, ``A simplex method for function minimization,'' \emph{The Computer Journal}, vol.~7, no.~4, pp.~308--313, 1965.
\bibitem{lagarias98} J.~C.~Lagarias, J.~A.~Reeds, M.~H.~Wright and P.~E.~Wright, ``Convergence properties of the Nelder--Mead simplex method in low dimensions,'' \emph{SIAM J.\ Optimization}, vol.~9, no.~1, pp.~112--147, 1998.
\bibitem{astrom98} K.~J.~\AA{}str\"om, H.~Panagopoulos and T.~H\"agglund, ``Design of PI controllers based on non-convex optimization,'' \emph{Automatica}, vol.~34, no.~5, pp.~585--601, 1998.
\bibitem{panagopoulos02} H.~Panagopoulos, K.~J.~\AA{}str\"om and T.~H\"agglund, ``Design of PID controllers based on constrained optimisation,'' \emph{IEE Proc.\ Control Theory Appl.}, vol.~149, no.~1, pp.~32--40, 2002.
\bibitem{garpinger08} O.~Garpinger and T.~H\"agglund, ``A software tool for robust PID design,'' \emph{17th IFAC World Congress}, Seoul, 2008.
\bibitem{lambrechts05} P.~Lambrechts, M.~Boerlage and M.~Steinbuch, ``Trajectory planning and feedforward design for electromechanical motion systems,'' \emph{Control Engineering Practice}, vol.~13, no.~2, pp.~145--157, 2005.
\bibitem{boerlage03} M.~Boerlage, M.~Steinbuch, P.~Lambrechts and M.~van de Wal, ``Model-based feedforward for motion systems,'' \emph{Proc.\ IEEE Conf.\ Control Applications}, pp.~1158--1163, 2003.
\bibitem{vandermeulen08} S.~H.~van der Meulen, R.~L.~Tousain and O.~H.~Bosgra, ``Fixed structure feedforward controller design exploiting iterative trials: application to a wafer stage and a desktop printer,'' \emph{J.\ Dynamic Systems, Measurement, and Control}, vol.~130, 051006, 2008.
\bibitem{armstrong94} B.~Armstrong-H\'elouvry, P.~Dupont and C.~Canudas de Wit, ``A survey of models, analysis tools and compensation methods for the control of machines with friction,'' \emph{Automatica}, vol.~30, no.~7, pp.~1083--1138, 1994.
\bibitem{ljung99} L.~Ljung, \emph{System Identification: Theory for the User}, 2nd ed., Prentice Hall, 1999.
\bibitem{soderstrom07} T.~S\"oderstr\"om, ``Errors-in-variables methods in system identification,'' \emph{Automatica}, vol.~43, no.~6, pp.~939--958, 2007.
\bibitem{koren80} Y.~Koren, ``Cross-coupled biaxial computer control for manufacturing systems,'' \emph{J.\ Dynamic Systems, Measurement, and Control}, vol.~102, no.~4, pp.~265--272, 1980.
\end{thebibliography}
\end{document}
